\chapter{MỞ ĐẦU}
\label{Chapter1}
\textit{Nội dung Chương 1 giới thiệu thực trạng chung của nhu cầu tra cứu cũng như đặt ra những câu hỏi liên quan đến vấn đề pháp luật ở Việt Nam, từ đó hình thành nên lý do và mục tiêu để thực hiện đồ án.}
\section{Giới thiệu chung}
	Hiện nay, nhu cầu tìm kiếm cũng như hỏi đáp các vấn đề liên quan đến các lĩnh vực trong đời sống đang rất lớn, đặc biệt là sử dụng các dịch vụ tra cứu trực tuyến. Các lĩnh vực có nhu cầu tìm kiếm thông tin như công nghệ, chính trị, xã hội, luật pháp,.... Đặc biệt là trong lĩnh vực luật pháp, điều mà trước đây chỉ có thể sử dụng phương pháp truyền thống là đến các cơ sở có khả năng pháp lý như văn phòng luật sư, các cơ sở có quyền về pháp lý,... mà người có nhu cầu muốn giải đáp thắc mắc có thể đến và đặt ra những câu hỏi cụ thể để từ đó, những chuyên gia có thể giải đáp được những câu hỏi cho họ.
	
	Nhưng với thời điểm công nghệ 4.0 hiện nay, việc sử dụng internet đối với mỗi người chúng ta là điều thiết yếu trong cuộc sống. Vậy nên, việc xuất hiện các trang web, những dịch vụ trực tuyến nhằm giải đáp những vấn đề trên được hình thành. Tuy nhiên, có một số điểm yếu mà còn tồn tại trong hệ thống này là:
	\begin{itemize}
		\item Các câu hỏi được giải đáp bởi các chuyên gia và chỉ được lưu trữ dạng văn bản thuần.
		\item Khi cần tra cứu câu hỏi, người dùng cần phải nhập đúng chính xác câu hỏi hoặc chính xác từ khoá mà có câu hỏi mẫu đã được lưu trữ trong cơ sở dữ liệu của hệ thống.
		\item Nếu như người dùng cần tìm những câu hỏi có nội dung tương tự như câu hỏi đang cần hỏi thì cần phải biết được lĩnh vực mà câu hỏi đã được hỏi.
	\end{itemize}


\section{Mục tiêu đề tài}
Với nhu cầu trao đổi thông tin hiện nay ngày càng cao cùng với sự phát triển của mạng Internet và công nghệ, hằng ngày con người luôn phải xử lý một lượng thông tin lớn. Do đó việc hỏi đáp mang tính thiết yếu hơn trong cuộc sống, và nhiều trang web, hệ thống hỏi đáp đã được xây dựng nên để đáp ứng nhu cầu đó. Nhưng trong số đó, vẫn còn nhiều hệ thống hỏi đáp chưa thể đưa ra câu trả lời ngay lập tức được hoặc phải do chính nhân tố con người trả lời một cách thủ công. Đặc biệt khi nói về lĩnh vực Pháp luật – lĩnh vực mang tính chất khô khan và cần nhiều sự chính xác và nhanh chóng khi đưa ra câu trả lời đối với một hệ thống hỏi đáp thì hiện nay vẫn chưa có nhiều hệ thống đáp ứng được.

Đề tài hướng đến mục tiêu là xây dựng một hệ thống hỏi đáp tự động tiếng Việt cho một lĩnh vực cụ thể là về pháp luật Việt Nam. Bằng cách áp dụng công nghệ nhiều hơn, ở đây chính là công nghệ trí tuệ nhân tạo giúp nâng cao hiệu quả, tốc độ và sự chính xác trong câu trả lời mà hệ thống đưa ra. Với hướng tiếp cận xử lý ngôn ngữ tự nhiên, sử dụng các Pre-trained Model (mô hình tiền huấn luyện) để phân tích ngữ nghĩa của các câu hỏi và cho ra kết quả chính xác nhất thì việc xây dựng một hệ thống về hỏi đáp Pháp luật Việt Nam là một đề tài hoàn toàn khả thi.

\pagebreak
\section{Nội dung của đồ án}

Đồ án có tổng cộng 7 chương với nội dung chính như sau:

\textbf{Chương 1:} Mở đầu

    Nội dung chương 1 giới thiệu về thực trạng về nhu cầu trong thực tế của việc hỏi đáp các vấn đề liên quan đến luật pháp ở Việt Nam, đưa ra các điểm yếu trong phương pháp truyền thống và đưa ra mục tiêu, phương pháp áp dụng của đồ án vào trong ứng dụng.

\textbf{Chương 2:} Khảo sát hiện trạng
	Nội dung chương 2 là báo cáo từ một khảo sát nhỏ của nhóm nhằm xác định được rằng trong thực tế thì có những nhóm người dùng nào sẽ có nhu cầu tìm kiếm, hỏi đáp về pháp luật và cách thức để người dùng tìm được đáp án cho vấn đề đang cần tìm kiếm.

\textbf{Chương 3:} Công nghệ sử dụng

    Nội dung của chương 2 là đưa ra các cơ sở lý thuyết để áp dụng vào trong đề tài, chỉ ra các công nghệ được áp dụng và phối hợp với nhau để giải quyết được vấn đề cần phải giải quyết trong đề tài nghiên cứu.

\textbf{Chương 4:} Thu Thập Dữ Liệu

    Nội dung chương 3 là trình bày về tổng quan của hệ thống và quy trình để xây dựng nên hệ thống về phần Back-end và Front-end.

\textbf{Chương 5:} Kiến trúc phần Back-end
	\begin{itemize}
		\item Trình bày chi tiết về quá trình thu thập dữ liệu pháp luật
		\item Phương pháp xây dựng cơ sở dữ liệu và lưu trữ dữ liệu
		\item Mô hình đào tạo dữ liệu phương pháp phân tích đưa ra được thông tin chuẩn xác từ mô hình. 
	\end{itemize}

\textbf{Chương 6:} Kiến trúc phần Front-end

    Trình bày về trình tự và cách để đưa dữ liệu được trả về ở phần Back-end hiển thị lên giao diện người dùng.

\textbf{Chương 7:} Kết luận và hướng phát triển
    Nội dung chương 5 là tóm tắt lại kết quả đạt được sau khi hoàn thành đề tài, ý nghĩa của đề tài và định hướng phát triển của đề tài trong tương lai
